\documentclass[11pt,a4paper]{article}

% ---------------------------------------------------------
% PREAMBLE (stable, warning-free, Overleaf-safe)
% ---------------------------------------------------------
\usepackage[utf8]{inputenc}
\usepackage[T1]{fontenc}
\usepackage{lmodern}
\usepackage{amsmath, amssymb, amsfonts}
\usepackage{bm}
\usepackage{geometry}
\usepackage{graphicx}
\usepackage{hyperref}
\usepackage{cite}
\usepackage{authblk}
\usepackage{physics}
\usepackage{mathtools}
\usepackage{textgreek}
\geometry{margin=2.4cm}

\title{\textbf{Companion LXXIV: Fisher Information Geometry of Persistence Diagrams in the Relaxation-Driven Cyclic Cosmology}}
\author{Michael Lehmann \\ \texttt{mi.lehmann@gmx.de}}
\date{April 2026}

\begin{document}
\maketitle

\begin{abstract}
Fisher information geometry provides a Riemannian structure on statistical 
manifolds, quantifying the sensitivity of probability distributions to underlying 
model parameters. Applied to persistence diagrams, this framework measures how 
topological features respond to changes in cosmological dynamics. Building on the 
Hilbert-space embeddings, kernel methods, and statistical tests developed in 
previous Companions, this work introduces the Fisher information metric for 
persistence diagrams in the Relaxation-Driven Cyclic Cosmology (RDCC). The 
relaxation mechanism modifies the scale-dependent evolution of the gravitational 
potential, inducing characteristic changes in the Fisher geometry of topological 
ensembles. The resulting information metric provides a powerful tool for 
parameter sensitivity, model discrimination, and information-theoretic analysis of 
weak-lensing topology.
\end{abstract}
% ---------------------------------------------------------
\section{Introduction}

Persistence diagrams encode the birth and death of topological features in weak-
lensing convergence fields. Previous Companions developed Hilbert-space embeddings, 
kernel PCA, and maximum mean discrepancy (MMD) as tools for analyzing and 
distinguishing the Relaxation-Driven Cyclic Cosmology (RDCC) from ΛCDM.

Fisher information geometry extends this framework by introducing a Riemannian 
metric on the space of probability distributions of persistence diagrams. This 
metric quantifies how sensitively the distribution responds to changes in 
cosmological parameters. In RDCC, the relaxation mechanism modifies the 
gravitational potential in a scale-dependent manner, producing characteristic 
signatures in the Fisher information of topological ensembles.
% ---------------------------------------------------------
\section{Fisher Information Geometry}

Given a family of probability densities $p(x|\theta)$ depending on parameters 
$\theta = (\theta_{1}, \dots, \theta_{d})$, the Fisher information matrix is

\begin{equation}
    I_{ab}(\theta)
    = \mathbb{E}
      \left[
        \partial_{\theta_{a}} \log p(x|\theta)
        \,
        \partial_{\theta_{b}} \log p(x|\theta)
      \right].
\end{equation}

This defines a Riemannian metric on the parameter manifold.

For persistence diagrams, we consider the induced density in an RKHS via the 
embedding $\Phi(D)$:



\[
    p(D|\theta)
    \propto
    \exp\left(
        - \| \Phi(D) - \mu(\theta) \|_{\mathcal{H}}^{2}
    \right),
\]



where $\mu(\theta)$ is the mean embedding.
% ---------------------------------------------------------
\section{RDCC Modification of Persistence Diagram Distributions}

In RDCC, the convergence evolves as

\begin{equation}
    \kappa_{\mathrm{RDCC}}(\ell)
    = \kappa(\ell)
      \left[
        1 - \chi_{\kappa}
        \left(\frac{\ell}{\ell_{\mathrm{rel}}}\right)^{\alpha}
      \right].
\end{equation}

This modifies the distribution of persistence diagrams:



\[
    p_{\mathrm{RDCC}}(D|\theta)
    = p(D|\theta)
      \left[
        1 - \chi_{p}
        \left(\frac{\ell}{\ell_{\mathrm{rel}}}\right)^{\alpha}
      \right].
\]



The mean embedding shifts as



\[
    \mu_{\mathrm{RDCC}}(\theta)
    = \mu(\theta)
      \left[
        1 - \chi_{\mu}
        \left(\frac{\ell}{\ell_{\mathrm{rel}}}\right)^{\alpha}
      \right].
\]


% ---------------------------------------------------------
\section{Fisher Information in RDCC}

The RDCC Fisher information matrix becomes

\begin{equation}
    I^{\mathrm{RDCC}}_{ab}
    = I_{ab}
      \left[
        1 - \chi_{I,ab}
        \left(\frac{\ell}{\ell_{\mathrm{rel}}}\right)^{\alpha}
      \right].
\end{equation}

RDCC induces:

\begin{itemize}
    \item enhanced sensitivity to large-scale topological modes,
    \item suppressed sensitivity to small-scale modes,
    \item a characteristic turnover near $\ell_{\mathrm{rel}}$,
    \item modified curvature of the information manifold,
    \item altered geodesic distances between topological distributions.
\end{itemize}
% ---------------------------------------------------------
\section{Information-Geometric Signatures of RDCC}

The Fisher metric reveals how RDCC modifies the information content of topological 
features. Key signatures include:

\begin{itemize}
    \item increased Fisher curvature for long-lived persistence features,
    \item reduced curvature for short-lived features,
    \item anisotropic deformation of the parameter manifold,
    \item enhanced distinguishability of RDCC from ΛCDM at large scales,
    \item modified geodesic flow of topological distributions across redshift.
\end{itemize}

These effects provide a direct information-theoretic probe of the RDCC 
gravitational field.
% ---------------------------------------------------------
\section{Conclusion}

Fisher information geometry provides a powerful framework for analyzing the 
parameter sensitivity of persistence diagrams in the Relaxation-Driven Cyclic 
Cosmology. The relaxation mechanism modifies the Fisher metric in a scale-dependent 
manner, producing distinctive signatures in the information geometry of weak-
lensing topology. These effects offer a direct, information-theoretic window into 
the relaxation scale and the temporal structure of the RDCC gravitational field. 
The present Companion establishes the foundation for future studies of Fisher 
kernels, parameter forecasts, and information geometry in cosmological topology.

\bibliographystyle{apsrev4-2}
\nocite{*}
\bibliography{references}


\end{document}
